\documentclass[10pt]{article}
\usepackage[T1]{fontenc}
\usepackage[utf8]{inputenc}
\usepackage[french]{babel}
\usepackage{amsmath,amssymb,mathrsfs,mathpazo}
\usepackage{geometry}
\geometry{a4paper,margin=2cm}
\usepackage{graphicx,array,tabularx,tabularray,colortbl,multirow,multicol,booktabs,enumitem,pifont,ulem,fancybox,diagbox,cancel}
\usepackage{pstricks,pstricks-add,pst-plot,pst-tree,pst-node,pst-text,pst-func,pst-math,pst-bspline,pst-3dplot,pst-eucl,pst-solides3d}
\let\pscalculate\undefined
\usepackage{pst-bezier}
\usepackage{tikz,pgfplots,tkz-tab}
\usetikzlibrary{arrows,arrows.meta,patterns,calc,positioning,trees,decorations.markings,shapes.geometric,shapes.misc,shapes.arrows,intersections,decorations.pathmorphing,decorations.pathreplacing,backgrounds,shadows,babel}
\pgfplotsset{compat=1.18}
\usepgfplotslibrary{fillbetween}
\usepackage[np]{numprint}
\usepackage{listings,makecell,pas-tableur,twoopt}
\let\barre\undefined
\usepackage[pstricks]{tablvar}
\definecolor{darkgreen}{RGB}{0,100,0}
\providecommand{\red}{\color{red}}
\providecommand{\blue}{\color{blue}}
\providecommand{\green}{\color{green}}
\providecommand{\black}{\color{black}}
\providecommand{\white}{\color{white}}
\providecommand{\gray}{\color{gray}}
\providecommand{\cyan}{\color{cyan}}
\providecommand{\magenta}{\color{magenta}}
\providecommand{\yellow}{\color{yellow}}
\providecommand{\violet}{\color{violet}}
\providecommand{\orange}{\color{orange}}
\providecommand{\euro}{\texteuro}
\providecommand{\EUR}{\texteuro}
\providecommand{\ieme}{\textsuperscript{e}}
\providecommand{\degres}{\ensuremath{^\circ}}
\providecommand{\degree}{\ensuremath{^\circ}}
\providecommand{\vv}[1]{\overrightarrow{#1}}
\providecommand{\Coord}[1]{#1}
\providecommand{\decosix}{}
\providecommand{\starredbullet}{\ensuremath{\bullet}}
\pagestyle{empty}
\setlength{\parindent}{0pt}
\usepackage{adjustbox}\DeclareUnicodeCharacter{2236}{:}\DeclareUnicodeCharacter{2190}{\ensuremath{\leftarrow}}
\providecommand{\R}{\mathbb{R}}
\providecommand{\N}{\mathbb{N}}
\providecommand{\Z}{\mathbb{Z}}
\providecommand{\Q}{\mathbb{Q}}
\providecommand{\D}{\mathbb{D}}
\providecommand{\C}{\mathbb{C}}
\providecommand{\e}{{\mathrm{e}}}
\providecommand{\vect}[1]{\overrightarrow{#1}}
\providecommand{\vectt}[1]{\overrightarrow{\mathrm{#1}}}
\providecommand{\vvt}[1]{\overrightarrow{\mathrm{#1}}}
\providecommand{\np}[2][]{{#2~#1}}
\providecommand{\vv}[1]{{\overrightarrow{#1}}}
\renewcommand{\d}{{\mathrm{d}}}
\providecommand{\degres}{{^\circ}}
\providecommand{\degre}{{^\circ}}
\providecommand{\blue}{\color{blue}}
\providecommand{\red}{\color{red}}
\providecommand{\enskip}{\quad}
\providecommand{\llbracket}{\lbrack\!\lbrack}
\providecommand{\rrbracket}{\rbrack\!\rbrack}
\providecommand{\euro}{{\text{€}}}
\providecommand{\no}{{\text{n°}}}
\providecommand{\interval}[2]{{[#1;#2]}}
\providecommand{\intervFO}[2]{{[#1;#2[}}
\providecommand{\dotfill}{\ldots}

\providecommand{\ds}{\displaystyle}
\providecommand{\textding}[1]{#1}
\providecommand{\Oij}{\(\mathrm{O};\vec{\imath},\vec{\jmath}\)}
\providecommand{\Oijk}{\(\mathrm{O};\vec{\imath},\vec{\jmath},\vec{k}\)}
\providecommand{\Ouv}{\(\mathrm{O};\vec{u},\vec{v}\)}
\providecommand{\intervFF}[2]{\ensuremath{[#1;#2]}}
\renewcommand{\euro}{\ifmmode\text{\texteuro}\else\texteuro\fi}
\DeclareUnicodeCharacter{00B0}{\ifmmode^{\circ}\else\textdegree\fi}
\DeclareUnicodeCharacter{20AC}{\ifmmode\text{\texteuro}\else\texteuro\fi}
% Énoncés : Éducation nationale. Transcription et corrigés : APMEP, Denis Vergès et contributeurs. Découpage et classement : CoopMaths, Sébastien Lozano, Éric Elter et contributeurs. Adaptation de la présentation : Maths Explorer.
% CC BY-SA : https://coopmaths.fr/www/about/
% https://coopmaths.fr/alea/static/e3c/2020/tex/e3c_2020_00_sujet12_4_cor.tex
\begin{document}
\noindent\textit{APMEP / CoopMaths. CC BY-SA. Présentation adaptée par Maths Explorer.}\par\medskip
\subsection*{1. On note $X$ la variable aléatoire qui prend pour valeur le nombre total de buts marqués à l'issue de cette série de tirs par Karim.}
	
	\subsubsection*{a. Réaliser un arbre pondéré permettant de décrire toutes les issues possibles.}
	\begin{center}
		\begin{tikzpicture}
			[level 1/.style={level distance=2cm,
				sibling distance=4cm},
			level 2/.style={level distance=2.5cm,
				sibling distance=2cm},
			level 3/.style={level distance=2.5cm,
					sibling distance=1cm}]
			\node {} [grow'=right]
			child {node {$M$}
				child {node {$M$}
						child {node {$M$}
						edge from parent node[above] {$0,7$}
					}
					child {node {$R$}
						edge from parent node[below] {$0,3$}
					}
				edge from parent node[above] {$0,7$}
				}
				child {node {$R$}
						child {node {$M$}
						edge from parent node[above] {$0,7$}
					}
					child {node {$R$}
						edge from parent node[below] {$0,3$}
					}
				edge from parent node[below] {$0,3$}
				}
				edge from parent node[above] {$0,7$}
			}
			child {node {$R$}
				child {node {$M$}
						child {node {$M$}
						edge from parent node[above] {$0,7$}
					}
					child {node {$R$}
						edge from parent node[below] {$0,3$}
					}
					edge from parent node[above] {$0,7$}
				}
				child {node {$R$}
						child {node {$M$}
						edge from parent node[above] {$0,7$}
					}
					child {node {$R$}
						edge from parent node[below] {$0,3$}
					}
					edge from parent node[below] {$0,3$}
				}
				edge from parent node[below] {$0,3$}
			}
			;
		\end{tikzpicture}
	\end{center}

	
	\subsubsection*{b. Déterminer la loi de probabilité de $X$.}
	
	$X$ peut prendre les valeurs 0, 1, 2 et 3.
	
	\[
	\begin{array}{|c|c|c|c|c|}
		\hline
			X & 0 & 1 & 2 & 3 \\
		\hline
		P(X = x_i) & 0,027 & 0,189 & 0,441 & 0,343 \\
			\hline
	\end{array}
	\]
	
	\subsubsection*{c. Calculer l'espérance $\mathbb{E}(X)$ de la variable aléatoire $X$.}
	
	\[
	\mathbb{E}(X) = 0 \times 0,027 + 1 \times 0,189 + 2 \times 0,441 + 3 \times 0,343 = 0,189 + 0,882 + 1,029 = 2,100
	\]
	
	Sur un grand nombre de tirs, Karim en réussira 21 sur 30.
	
	\subsection*{2. On propose à un spectateur le jeu suivant : il mise 15 € avant la série de tirs au but de Karim ; chaque but marqué par Karim lui rapporte 6 €, et chaque but manqué par Karim ne lui rapporte rien.}
	
	\subsubsection*{a. Exprimer $Y$ en fonction de $X$.}
	
	On a $Y = 6X - 15$. Les valeurs de $Y$ sont donc : -15, -9, -3 et 3.
	
	\subsubsection*{b. Calculer l'espérance $\mathbb{E}(Y)$ de la variable aléatoire $Y$. Interpréter ce résultat dans le contexte de l'énoncé.}
	
	\[
	\mathbb{E}(Y) = -15 \times 0,027 - 9 \times 0,189 - 3 \times 0,441 + 3 \times 0,343 = -0,405 - 1,701 + 1,029 = -2,4
	\]
	
	Cela signifie qu'en moyenne sur un grand nombre de tirs, le spectateur perdra 2,40 € par série de trois tirs. Le jeu est inéquitable.

\end{document}
